\chapter*{Conclusion générale}
\addcontentsline{toc}{chapter}{Conclusion générale}

Ce projet a étudié la faisabilité d'un système de détection précoce des AVC à partir de données fortement déséquilibrées. L'exploration a notamment mis en évidence une association importante entre l'absence de mesure de l'IMC et la prévalence observée de l'AVC : celle-ci atteint 19,90\,\% lorsque l'IMC est manquant, contre 4,26\,\% lorsqu'il est renseigné. Ce constat a motivé la création de l'indicateur \texttt{bmi\_missing}, sans interpréter cette association comme une relation causale ni conclure à lui seul sur le mécanisme de non-réponse.

Le travail de préparation a conservé cette information, plafonné les valeurs extrêmes de l'IMC au lieu d'éliminer les dossiers concernés et séparé le jeu train/test avant l'ajustement des statistiques d'imputation et de standardisation. Les opérations apprises ont ainsi été ajustées sur le train, puis appliquées au test. Le benchmark a comparé plusieurs familles de modèles et stratégies de pondération ou de rééchantillonnage ; l'encapsulation de SMOTE dans un pipeline a permis de limiter son application aux plis d'entraînement pendant la validation croisée. Les modèles arborescents ont par ailleurs été régularisés afin de réduire leur complexité.

L'apport méthodologique central est l'optimisation du seuil à partir de probabilités \textit{out-of-fold} calculées sur le train. Le critère \(F_2\) a privilégié le rappel, tout en maintenant un plancher de spécificité fixé à 0,70 pour encadrer la charge de faux positifs. Ce plancher constitue une contrainte opérationnelle du projet, et non une norme clinique universelle. La règle de sélection, établie sans consulter le test, a retenu la Random Forest. Son seuil est \(0.5792571442648802\) ; sur le test indépendant, elle obtient un rappel de 0,80 et une spécificité de 0,766, avec 10 faux négatifs parmi les 50 cas AVC. La précision de 0,150 et les 227 faux positifs observés montrent cependant que le compromis reste exigeant et doit être évalué avec les équipes cliniques.

Les résultats ne constituent pas une validation médicale ni une preuve de généralisation à d'autres établissements ou populations. Ils dépendent du jeu disponible, et le nombre de cas AVC dans le test reste limité. En outre, dans la version étudiée, les colonnes issues du One-Hot Encoding ont été établies avant le split ; une évolution devra ajuster et sérialiser l'encodeur exclusivement sur le train pour renforcer l'isolation du schéma des catégories. Des validations externes, une analyse de calibration et une étude des performances par sous-groupes sont également nécessaires avant tout usage opérationnel.

\section*{Perspectives}
La prochaine étape consiste à encapsuler le modèle exporté et ses artefacts de prétraitement dans un microservice FastAPI. Des schémas Pydantic pourront valider les champs, les types, les valeurs admissibles et les règles de cohérence des requêtes. Le service devra respecter l'ordre des caractéristiques et le seuil exact conservé dans \texttt{model\_metadata.json}, puis retourner un score et une catégorie d'aide au triage sans se substituer au jugement médical.

Ce microservice pourra ensuite servir d'interface backend à une application Next.js destinée aux urgentistes. La mise en production devra être précédée de tests d'intégration et de sécurité, d'une validation clinique appropriée, ainsi que d'un suivi continu de la calibration, des faux négatifs, de la charge d'alertes et de toute dérive des données.
